% Source: https://artofproblemsolving.com/community/c3h3267911_2023_lmt_fall_guts_round_p1p15_lexington_mathematical_tournament
% Stable identity: aops-topic-3267911-q12
For one middle digit, the alternating sum is $a-b+c$. For two middle digits it is $a-c$. Thus $a=c$ is necessary, because $a-c$ lies between-8 and 9 and must be divisible by 11. Then $2a-b\equiv0\pmod{11}$. Conversely these conditions ensure that any odd number of middle digits has alternating sum $2a-b$ and any even number has sum 0. For $a=1,2,3,4$, choose $b=2a$; for $a=6,7,8,9$, choose $b=2a-11$. The value $a=5$ would require $b=10$ or-1, impossible. Hence there are $\boxed8$ numbers:121,242,363,484,616,737,858,979.
